\subsection{Endpoints, relative paths, and the component theorem}

A word is read from its initial vertex to its terminal vertex, in the positional order of its vertices. A negative sign means that its generating arrow points against this positional order. It does not mean that the category $M$ or its arrows have been replaced by an opposite category. We concatenate the word for $a$ to $b$ first and the word for $b$ to $c$ second, while writing the composite path as $d\circ c$.

The definition at VII ms.~65--67 calls the relative fixed-type category a ``sous-cat\'egorie pleine''. Thus it restricts the backward arrows of objects, and does not restrict interior components of natural transformations. The text at ms.~66 saying that backward arrows ``sont dans $X$'' must read ``sont dans $\Sigma$'', as the supplied transcription already annotates. We use the precise two category constructions in B.1 throughout; the weak-vertical version is explicitly distinguished by its second superscript.

\begin{result}{A.P1}{Proved}{The localization theorem on components}
For a small category $M$, a wide subcategory $\Sigma$, and objects $a,b$, there is a natural bijection
\[
 \pi_0\operatorname{Ch}_\Sigma(M;a,b)
       \cong\operatorname{Hom}_{M[\Sigma^{-1}]}(a,b).
\]
It preserves concatenation and the empty-path identities. A backward edge is interpreted as the inverse of its arrow in the ordinary localization. No two-out-of-three, saturation, or model-category axiom is assumed.
\end{result}
\begin{proof}
Strict concatenation of the path categories defines a category $\Pi_\Sigma(M)$ with objects those of $M$ and the displayed component sets as its Hom-sets. Sending an arrow to its one-edge positive path gives a functor $M\to\Pi_\Sigma(M)$. For composable arrows, the type map $(+)\to(+,+)$ gives a path morphism from their composite to their concatenation. For an identity arrow, collapse the single edge to the empty type. These observations prove the functor equations on components.

If $w:a\to b$ belongs to $\Sigma$, its reversed one-edge path is allowed. The two concatenations with that reversed path have types $(+,-)$ and $(-,+)$. In the first, the constant path maps to the backtrack by a transformation with middle component $w$; in the second, the backtrack maps to the constant path by such a transformation. Type collapse identifies each constant path with the empty path. Thus the image of $w$ is invertible.

Let $F:M\to D$ be any functor sending $\Sigma$ to isomorphisms. For a path $c$, define its effect in $D$ by applying $F$, inverting the backward arrows, and composing. A type map does not change this effect. For a natural transformation of paths, its commuting squares give equal total effects because the endpoint components are identities. In checking a backward square one cancels the two invertible horizontal arrows; no vertical arrow is inverted. Hence the effect descends to a functor $\Pi_\Sigma(M)\to D$ extending $F$. It is unique because every path is the concatenation of its one-edge paths, and the backward ones have just been shown to be the prescribed inverses. This is the universal property of ordinary localization, with its structural functor from $M$ retained. It proves the bijection and all its stated compatibilities. If a Hom-set is empty, both sides are empty by the same path description; disconnected categories present no additional condition.
\end{proof}

\begin{result}{A.P2}{Proved}{The sequential models and the comparison roof}
Let $t:\mathbb N^{\mathrm{op}}\to\mathrm{Ty}$ be a projective type sequence. With $F_\tau$ as in B.1, define
\[
 C_\infty(t;a,b)=\operatorname*{colim}_{n\in\mathbb N}F_{t_n}(a,b)
\]
in $\mathrm{Cat}$, with transition maps given by precomposition with $t_{n+1}\to t_n$. Let $E_t$ be the fibred Grothendieck construction over $\mathbb N^{\mathrm{op}}$. There is a specified span
\[
 C_\infty(t;a,b)\ \longleftarrow\ E_t
       \ \longrightarrow\ \operatorname{Ch}_\Sigma(M;a,b).
\]
The left leg sends a path to its colimit class. The right leg retains its type and path and applies $t$ to a base arrow. The left leg is always in $W_\infty$. For every sequence cofinal in the reachability preorder, both legs are in $W_\infty$, by C.P3.
\end{result}
\begin{proof}
An arrow in $E_t$ from $(n,c)$ to $(m,d)$, where $n\geq m$, is an arrow $c\to(t_n\to t_m)^*d$ in $F_{t_n}$. In the ordinary colimit this is an arrow $[c]\to[d]$, giving the left leg. Its type map and transformation give the right leg. Identities and compositions agree with the formula in B.1. The left weak-equivalence assertion is the filtered homotopy-colimit theorem in A.2; no full-faithfulness of transitions is required. The right assertion retains the exact additional hypothesis and refers to the independent enriched-type proof C.P1--C.P3.

For the block $D=(+,-)$, the left and right one-sided systems are represented by the two half-fences, and the symmetrically padded system by the integer fence $J$ with edges $2k\to2k-1$ and $2k\to2k+1$. Paths are eventually constant on the indicated tails; fixed-endpoint transformations are eventually the specified endpoint identities. Constant-tail extension gives the actual $q_L,q_R$ of CATHOT E.23--E.25. For a general block containing both signs, the corresponding models are the ordinary colimits of powers with left, right, or symmetric outer collapse. C.P4 proves the stated padding comparisons and their compatibility with composition for those precise transitions. No universal-arrow lifting property of the endpoint map is asserted by the term ``path model''; Part E distinguishes it from smoothness.
\end{proof}
